\begin{tabular}{@{}l*{7}{r}@{}}
\toprule
\cabecalhoL{Ano} & \cabecalho{ICMS bruto\\R\$ bi} & \cabecalho{Outras\\deduções} & \cabecalho{Cota\\municipal} & \cabecalho{FECOP} & \cabecalho{Receita de\\referência} & \cabecalho{Fator de\\ajuste} & \cabecalho{Referência\\R\$ de 2025} \\
\midrule
2019 & 11,36 & 0,00 & 2,84 & 0,05 & 8,57 & 1,741 & 14,93 \\
2020 & 11,87 & 0,01 & 2,97 & 0,06 & 8,96 & 1,705 & 15,27 \\
2021 & 15,33 & 0,00 & 3,83 & 0,07 & 11,57 & 1,358 & 15,71 \\
2022 & 16,58 & 0,00 & 4,14 & 0,08 & 12,51 & 1,271 & 15,89 \\
2023 & 17,66 & 0,00 & 4,41 & 0,08 & 13,32 & 1,235 & 16,45 \\
2024 & 20,80 & 0,01 & 5,20 & 0,09 & 15,69 & 1,073 & 16,84 \\
2025 & 22,52 & 0,01 & 5,63 & 0,11 & 16,99 & 1,000 & 16,99 \\
2026 (est.) & 22,99 & 0,01 & 5,74 & 0,11 & 17,34 & --- & 17,34 \\
\midrule
\multicolumn{7}{@{}l}{\textbf{Média da receita de referência ajustada}} & \textbf{16,18} \\
\multicolumn{7}{@{}l}{\textbf{Arrecadação nacional de referência em 2025}} & \textbf{1.020,95} \\
\multicolumn{7}{@{}l}{\textbf{Coeficiente de participação do Espírito Santo}} & \textbf{1,5846\%} \\
\bottomrule
\end{tabular}
